% Upload this file to Overleaf and select pdfLaTeX as the compiler.
% No external .bib file, code files, or special compiler options are needed.
% Upload billing.png together with this source, or import the supplied ZIP.
\documentclass[10pt,letterpaper]{article}
\usepackage[margin=0.7in]{geometry}
\usepackage[T1]{fontenc}
\usepackage{lmodern}
\usepackage{microtype}
\DisableLigatures{encoding = *, family = tt*}
\usepackage{amsmath,amssymb}
\usepackage{booktabs,tabularx,array}
\usepackage{xcolor,graphicx,listings,enumitem,fancyhdr}
\usepackage{xurl}
\usepackage[colorlinks=true,linkcolor=blue!45!black,urlcolor=blue!45!black]{hyperref}

% ---------- Report details ----------
\newcommand{\BillingImage}{billing.png}
% ------------------------------------

\definecolor{codebg}{HTML}{F4F6F8}
\setlength{\parindent}{0pt}
\setlength{\parskip}{4pt}
\setlength{\emergencystretch}{2em}
\setlist{nosep,leftmargin=*}
\renewcommand{\arraystretch}{1.12}
\newcolumntype{Y}{>{\raggedright\arraybackslash}X}
\lstset{basicstyle=\ttfamily\footnotesize,backgroundcolor=\color{codebg},
  frame=single,rulecolor=\color{black!12},breaklines=true,
  columns=fullflexible,keepspaces=true,showstringspaces=false,upquote=true,
  aboveskip=7pt,belowskip=7pt,xleftmargin=4pt,xrightmargin=4pt}
\pagestyle{fancy}
\fancyhf{}
\fancyhead[L]{\small Homework -2}
\fancyhead[R]{\small U44607487}
\fancyfoot[C]{\small\thepage}
\setlength{\headheight}{14pt}
\hypersetup{pdftitle={Homework 2: Cloud Storage and Graph Analysis},
  pdfauthor={Srinivasa Sai Chava},pdfsubject={Cloud Storage and Graph Analysis}}
\newcommand{\repo}{\url{https://github.com/Chava-Sai/Cloud-Computing.git}}

\begin{document}
\thispagestyle{empty}
\noindent
\begin{minipage}[t]{0.59\linewidth}
\vspace{0pt}
{\LARGE\bfseries Homework -2}\par
\vspace{4pt}
{\large Cloud Storage and Graph Analysis}
\end{minipage}\hfill
\begin{minipage}[t]{0.39\linewidth}
\vspace{0pt}
\raggedleft
{\fontsize{14}{18}\selectfont\bfseries Srinivasa Sai Chava\par}
\vspace{4pt}
{\fontsize{14}{18}\selectfont U44607487\par}
\end{minipage}
\par\vspace{12pt}

\section{Project and dataset}
The program reads 12,000 HTML pages from Google Cloud Storage, constructs a
directed graph, and computes link statistics, PageRank, and outward closeness
centrality. Graph parsing and algorithms are implemented using Python's standard
library; no graph library is used. Only downloads use parallel workers, through
Google's storage transfer manager, as permitted by the instructor's Piazza
clarification. Parsing and graph computations run on one thread.

\begin{tabularx}{\linewidth}{@{}lY@{}}
\toprule
Project ID & \texttt{thermal-circle-508221-m7} \\
Bucket & \texttt{thermal-circle-508221-m7-hw2} \\
Location / storage class & \texttt{us-central1} / Standard \\
Object prefixes & \texttt{hw2/} (original); \texttt{hw2-gzip/} (gzip encoding) \\
Dataset & \texttt{0.html} through \texttt{11999.html}; 1,946,934 link occurrences \\
Code repository & \repo \\
Measured source revision & \texttt{2dd9507d8905feb5878bf435c7c910dbff195396} \\
\bottomrule
\end{tabularx}

\subsection{Generating and uploading the pages}
The supplied generator's logic was retained, including \texttt{random.seed(0)}.
It was run with 12,000 files and the maximum-reference parameter set to 325:
\begin{lstlisting}
mkdir -p files_12k
cd files_12k
python3 ../generate.py -n 12000 -m 325
cd ..
\end{lstlisting}
The generator uses exclusive upper bounds in both \texttt{randrange(1,max\_refs)}
and \texttt{range(1,num\_refs)}. Consequently, its actual outgoing link count
ranges from 0 to 323 when the parameter is 325. Duplicate links and self-links
are part of the generated data and are preserved.

\subsection{Bucket setup and public permissions}
Billing was enabled for the project. The following commands created the regional
bucket, granted public reading and listing, and uploaded the original pages:
\begin{lstlisting}
gcloud config set project thermal-circle-508221-m7
gcloud storage buckets create gs://thermal-circle-508221-m7-hw2 \
  --project=thermal-circle-508221-m7 --location=us-central1 \
  --default-storage-class=STANDARD --uniform-bucket-level-access
gcloud storage buckets add-iam-policy-binding \
  gs://thermal-circle-508221-m7-hw2 \
  --member=allUsers --role=roles/storage.objectViewer
gcloud storage rsync files_12k \
  gs://thermal-circle-508221-m7-hw2/hw2/ --recursive \
  --project=thermal-circle-508221-m7
\end{lstlisting}
Uniform bucket-level access uses IAM instead of per-object ACLs. The
\texttt{allUsers} binding grants anonymous object reading and listing. The
\texttt{hw2/} directory is an object-name prefix; no separate directory object
is required. The bucket remains public for grading.

\clearpage
\section{Running and verifying the program}
The public repository can be cloned without repository credentials. The Python
program uses an anonymous storage client by default, so graders do not need
Google credentials to read the bucket.
\begin{lstlisting}
git clone https://github.com/Chava-Sai/Cloud-Computing.git
cd Cloud-Computing/HW-2
python3 -m venv .venv
source .venv/bin/activate
python3 -m pip install -r requirements.txt
python3 -m unittest discover -v
python3 -u analyze.py \
  --bucket thermal-circle-508221-m7-hw2 \
  --prefix hw2/ --expected-nodes 12000 \
  --download-workers 32 --convergence both --verbose-pr
\end{lstlisting}
For the completed Cloud Shell run, the same command used
\texttt{--prefix hw2-gzip/}. That prefix is also available to graders.
To reproduce the recorded dependency versions, install
\texttt{results/laptop-packages.txt} instead of \texttt{requirements.txt}.
The measurements used Python 3.13.7 on the laptop, 3.12.3 in Cloud Shell, and
3.11.2 on the VM. Commands are run from the repository's \texttt{HW-2} folder.

\subsection{VM setup}
An \texttt{e2-medium} VM was created in \texttt{us-central1-a}, in the same
region as the bucket, with Debian 12 and a 10 GB standard persistent boot disk.
\begin{lstlisting}
gcloud services enable compute.googleapis.com \
  --project=thermal-circle-508221-m7
gcloud compute instances create hw2-vm \
  --project=thermal-circle-508221-m7 --zone=us-central1-a \
  --machine-type=e2-medium \
  --image-family=debian-12 --image-project=debian-cloud \
  --boot-disk-size=10GB --boot-disk-type=pd-standard \
  --no-service-account --no-scopes
gcloud compute ssh hw2-vm \
  --project=thermal-circle-508221-m7 --zone=us-central1-a
\end{lstlisting}
Inside the VM, Python environment tools and Git were installed before cloning
the repository and running the commands above:
\begin{lstlisting}
sudo apt-get update
sudo apt-get install -y python3-pip python3-venv git
\end{lstlisting}
The VM did not need a service account because bucket access was public. An SSH
key was required for login; the working assignment key was selected with
\texttt{--ssh-key-file="\$HOME/.ssh/hw2\_compute\_engine"}.

\subsection{Checking public access}
The anonymous object-list request succeeded, and downloading \texttt{0.html}
returned HTTP 200 and 97,586 bytes. These commands repeat those checks:
\begin{lstlisting}
curl --fail --silent --show-error \
 'https://storage.googleapis.com/storage/v1/b/thermal-circle-508221-m7-hw2/o?prefix=hw2%2F&maxResults=1'
curl --fail --silent --show-error --output /tmp/hw2-public-0.html \
 --write-out 'HTTP status: %{http_code}\n' \
 'https://storage.googleapis.com/thermal-circle-508221-m7-hw2/hw2/0.html'
\end{lstlisting}

\clearpage
\section{Graph representation and PageRank}
\subsection{Parsing and degree statistics}
Each numbered file is one node. A case-insensitive regular expression extracts
the generator's numeric \texttt{href="ID.html"} links. The parser targets the
supplied file format rather than arbitrary web HTML. Page IDs must be contiguous,
and links to IDs outside the dataset are rejected.

\texttt{out\_links[i]} stores every target occurrence in page \(i\).
The outgoing count is the list length. Reverse adjacency stores pairs of source
ID and multiplicity; incoming counts sum all link occurrences. A repeated link
therefore contributes more than once to degrees and PageRank. Both sums of
degrees equal the total number of link occurrences.

For incoming and outgoing counts separately, the program reports average,
median, maximum, minimum, and five quintile upper endpoints: P20, P40, P60,
P80, and P100. This gives nine values in each direction, following the Piazza
clarification. P100 repeats the maximum. For sorted values \(x_0,\ldots,x_{n-1}\),
the percentile at fraction \(q\) uses position \(h=(n-1)q\), interpolating
linearly between its neighboring indices.

\subsection{Iterative PageRank}
Every node starts with rank \(1/n\). Let \(d=0.85\), \(C(T)\) be the outgoing
link count, and \(m_{TA}\) the number of links from \(T\) to \(A\). Each iteration
uses only the previous rank vector:
\[
p_{k+1}(A)=\frac{0.15}{n}
 +0.85\sum_{T:C(T)>0}\frac{m_{TA}\,p_k(T)}{C(T)}
 +\frac{0.85}{n}\sum_{T:C(T)=0}p_k(T).
\]
The middle term is the assignment's incoming-link contribution formula. The
last term handles dangling pages (zero outgoing links) by distributing their
rank uniformly. This avoids division by zero and preserves total rank.

\subsection{Stopping rule and its limitation}
The assignment's literal stopping measurement is
\[
\Delta_{\mathrm{sum}}=100\,
\frac{|\sum_A p_{k+1}(A)-\sum_A p_k(A)|}{\sum_A p_k(A)}.
\]
It is available through \texttt{--convergence sum}, with the required default
threshold of 0.5\%. Since dangling redistribution conserves total rank, this
criterion can stop after one iteration even while individual ranks change.
The recorded runs use \texttt{--convergence both}, which additionally requires
\[
\Delta_{L1}=100\,
\frac{\sum_A|p_{k+1}(A)-p_k(A)|}{\sum_A p_{k+1}(A)}\leq0.5.
\]
Thus, both the specified total-rank condition and a rank-vector stability check
must hold. This additional check is an explicit implementation choice, not a
different interpretation attributed to the instructor. The program raises an
error if its iteration limit is reached without convergence. The reported
scores are approximate values at the selected tolerance.

\subsection{Cost of the computation}
With \(n\) nodes and \(m\) link occurrences, building the adjacency structure
requires \(O(n+m)\) work and storage, excluding the HTML text itself. Each
PageRank iteration traverses nodes and reverse edges, requiring at most
\(O(n+m)\) work. Degree percentiles require sorting the \(n\) counts.

\clearpage
\section{Closeness centrality and independent tests}
\subsection{Directed outward closeness}
Distances follow outgoing links. For a source \(u\), let \(R_u\) be the number
of other reachable nodes and \(S_u\) the sum of their shortest-path distances.
The score uses reachability normalization:
\[
\operatorname{CC}(u)=
\begin{cases}
\displaystyle\frac{R_u}{n-1}\frac{R_u}{S_u},&R_u>0\text{ and }n>1,\\[5pt]
0,&\text{otherwise}.
\end{cases}
\]
For a node that reaches the entire graph, this reduces to \((n-1)/S_u\).
The reachability factor prevents a node reaching only a tiny component from
receiving the same score as a well-connected node with equally short distances.
Duplicate edges do not shorten a path, and self-links add no new reachable node.

Breadth-first search (BFS) computes distances from every source, and the node
with the largest score is selected. The queue implementation explores nodes
in distance order. For graphs larger than 2,000 nodes, the implementation uses
Python integer bitmasks for neighbor sets, visited nodes, and BFS frontiers.
At each level it unions the outgoing neighbor masks and removes visited nodes;
the number of newly reached nodes multiplied by the level contributes to the
distance sum. The bitmask implementation still runs on one thread. Conventional
all-source BFS has \(O(n(n+m))\) worst-case work; bit operations reduce Python
loop overhead, but large integer operations are not constant-time.

\subsection{PageRank test with an exact solution}
An independent three-node graph has edges
\(0\to1\), \(0\to2\), \(1\to2\), and \(2\to0\). With \(d=0.85\), its stationary
ranks satisfy
\[
x=0.05+0.85z,\qquad y=0.05+0.425x,\qquad
z=0.05+0.425x+0.85y.
\]
Solving gives
\[
(x,y,z)=\left(\frac{686}{1769},\frac{380}{1769},\frac{703}{1769}\right).
\]
The regression test uses a tighter threshold of \(10^{-10}\) percent and checks
these expected fractions to 11 decimal places. Another exact test uses the
two-node graph \(0\to1\), where node 1 dangles; the expected ranks are
\((20/57,37/57)\). Symmetric cycles, duplicate links, rank conservation, and
early stopping under the literal sum rule are also tested.

\subsection{Closeness tests with independent distances}
For the directed path \(0\to1\to2\to3\), the scores are
\((1/2,4/9,1/3,0)\); node 0 is best. For a bidirectional four-node star,
the center's distances are \(1,1,1\), giving score 1; a leaf has distances
\(1,2,2\), giving score \(3/5\).
An exhaustive regression test constructs all \(2^9=512\) directed three-node
graphs, including self-links, and compares both BFS implementations against
an independently coded Floyd--Warshall distance algorithm. It also introduces
duplicate links. These tests do not use the randomly generated 12,000 pages.

\begin{tabularx}{\linewidth}{@{}Yrr@{}}
\toprule Environment & Tests passed & Test time (s) \\ \midrule
Laptop & 27 & 0.026 \\
Cloud Shell & 27 & 0.060 \\
VM & 27 & 0.072 \\ \bottomrule
\end{tabularx}
Other tests cover degree statistics, invalid or missing pages, pagination,
download failure handling, and completion of downloads before parsing begins.

\clearpage
\section{Analysis output}
All three successful runs loaded \textbf{12,000 nodes} and
\textbf{1,946,934 link occurrences}. They returned the same link statistics,
top five pages, and maximum closeness score. Incoming and outgoing averages
agree because every directed link contributes once to each total.

\subsection{Link statistics across all pages}
\begin{center}
\begin{tabular}{@{}lrr@{}}
\toprule Statistic & Incoming links & Outgoing links \\ \midrule
Average & 162.2445 & 162.2445 \\
Median & 162 & 162 \\
Maximum & 233 & 323 \\
Minimum & 115 & 0 \\
P20 & 152 & 65 \\
P40 & 159 & 130 \\
P60 & 165 & 195 \\
P80 & 173 & 260 \\
P100 & 233 & 323 \\
\bottomrule
\end{tabular}
\end{center}

\subsection{PageRank iterations and top five pages}
The \texttt{both} stopping rule completed after three iterations in every run:
\begin{center}
\begin{tabular}{@{}rrrr@{}}
\toprule Iteration & Total rank & Sum change (\%) & L1 change (\%) \\ \midrule
1 & 1.000000 & 0.000000 & 8.7758 \\
2 & 1.000000 & 0.000000 & 0.9758 \\
3 & 1.000000 & 0.000000 & 0.1158 \\
\bottomrule
\end{tabular}
\end{center}

\begin{center}
\begin{tabular}{@{}rlrrr@{}}
\toprule Rank & Page & PageRank & Incoming & Outgoing \\ \midrule
1 & \texttt{5207.html} & 0.00019553 & 183 & 144 \\
2 & \texttt{7443.html} & 0.00018286 & 196 & 86 \\
3 & \texttt{7400.html} & 0.00017964 & 159 & 247 \\
4 & \texttt{950.html} & 0.00017441 & 143 & 123 \\
5 & \texttt{369.html} & 0.00017332 & 185 & 7 \\
\bottomrule
\end{tabular}
\end{center}
PageRank is not just incoming degree: a link's contribution depends on the
source's rank and outgoing count. Therefore the page with the most incoming
links need not have the largest PageRank.

\subsection{Best closeness result}
\begin{lstlisting}
Best closeness centrality: page 10376.html  score = 0.50464735
\end{lstlisting}
The complete console logs, including progress messages and timings, are stored
in the repository under \texttt{HW-2/results/}:
\begin{itemize}
\item \texttt{laptop-parallel-bucket-run.txt}
\item \texttt{cloudshell-gzip-bucket-run.txt}
\item \texttt{vm-parallel-bucket-run.txt}
\end{itemize}
Each corresponding exit-code file contains \texttt{0}.

\clearpage
\section{Runtime comparison and Cloud Shell downloads}
Each completed run used 32 Google transfer-manager download workers, followed
by single-threaded parsing and graph calculations. The laptop used an Apple
M4 Pro with 24 GiB RAM and macOS 27.0. Cloud Shell reported an Intel Xeon at
2.20 GHz. The \texttt{e2-medium} VM reported an AMD EPYC 7B12. These are single
measurements rather than averages of repeated trials.

\begin{tabularx}{\linewidth}{@{}Yrrr@{}}
\toprule Stage (seconds) & Laptop & Cloud Shell & VM \\ \midrule
Object listing & 10.92 & 28.59 & 4.69 \\
Downloading & 111.23 & 136.99 & 112.08 \\
Parsing and graph construction & 1.23 & 4.02 & 3.50 \\
Combined loading stage & 123.99 & 170.32 & 120.99 \\
Link statistics & 0.001 & 0.005 & 0.004 \\
PageRank & 0.29 & 0.90 & 0.63 \\
Closeness & 102.67 & 255.93 & 339.36 \\
\textbf{Total wall time} & \textbf{226.96} & \textbf{427.16} & \textbf{460.99} \\
\bottomrule
\end{tabularx}
The combined loading stage already includes listing, downloading, parsing,
and other loader overhead; these component times must not be added twice.

\subsection{Observed Cloud Shell problem and resolution}
The laptop and VM used uncompressed objects under \texttt{hw2/}. In Cloud Shell,
the uncompressed 32-worker run had completed only 1,200 files after about
45 minutes, with low process CPU utilization. It later failed on
\texttt{hw2/2701.html} after read timeouts exhausted its 300-second retry budget.
The program reported failure rather than producing results from a partial graph.
Increasing concurrency alone did not resolve the slow transfers.

To reduce transferred bytes, the same generated files were uploaded under
\texttt{hw2-gzip/} with gzip content encoding:
\begin{lstlisting}
gcloud storage cp 'files_12k/*.html' \
  gs://thermal-circle-508221-m7-hw2/hw2-gzip/ \
  --gzip-local=html --project=thermal-circle-508221-m7
\end{lstlisting}
This preserves the numbered object names and original decoded HTML. A test of
\texttt{0.html} received 1,576 compressed bytes instead of 97,586 decoded bytes,
about a 62-fold reduction. The compressed request took 0.19 seconds; a normal
Python client request took 0.11 seconds and was verified equal to explicit gzip
decompression. The full Cloud Shell run then succeeded using the same program.

The low CPU usage during failed downloads, network read timeouts, and benefit
from smaller responses support a transfer bottleneck in that Cloud Shell
session. They do not identify a specific bandwidth limit or prove provider
throttling. The files' repeated filler text explains their high compressibility.

\subsection{Interpreting the measured differences}
Cloud Shell used gzip while the other completed runs used uncompressed objects;
the totals therefore are \textbf{not a controlled comparison of network speed}.
The graph and algorithms were the same. Laptop and VM download times were nearly
equal. Closeness was the main computation cost, and the laptop completed it
fastest. CPU performance, available CPU resources, and different Python versions
can contribute to this difference; the measurements do not isolate their effects.
An earlier sequential laptop baseline took 2,069.55 seconds overall, including
1,943.12 seconds loading. It used older download settings and is retained as a
separate baseline rather than mixed with the 32-worker measurements.

\clearpage
\section{Cleanup, billing, and AI disclosure}
\subsection{Deleting the VM}
Before deleting the VM, the completed output, test log, package versions,
environment details, and exit-code file were copied to the laptop. The saved
configuration confirms \texttt{e2-medium}, \texttt{us-central1-a}, Debian 12,
and the 10 GB boot disk. The following command deleted the VM and boot disk:
\begin{lstlisting}
gcloud compute instances delete hw2-vm \
  --project=thermal-circle-508221-m7 --zone=us-central1-a \
  --delete-disks=boot --quiet
\end{lstlisting}
The deletion command returned \texttt{Deleted} for \texttt{hw2-vm}.
Afterward, the following instance and disk queries both returned \texttt{[]}:
\begin{lstlisting}
gcloud compute instances list \
  --project=thermal-circle-508221-m7 \
  --filter="name=hw2-vm" --format=json
gcloud compute disks list \
  --project=thermal-circle-508221-m7 \
  --filter="name=hw2-vm" --format=json
\end{lstlisting}
The evidence is saved as \path{vm-deletion.txt},
\path{vm-after-deletion.json}, and \path{vm-disks-after-deletion.json}.
The bucket is retained so the instructor and teaching staff can run the code.

\subsection{AI use and understanding}
I used an AI coding assistant to help generate and revise the Python
implementation and tests, explain the graph algorithms, troubleshoot Google
Cloud setup and slow downloads, and prepare and format this report. The
assistance included code suggestions for parsing, PageRank, BFS-based closeness,
storage loading, and independent test cases. I ran the setup commands, tests,
and analysis on my laptop, Cloud Shell, and the VM, and collected the resulting
logs. The reported measurements are from those executions, not estimated or
generated timing results.

The main concepts behind the solution are:
\begin{itemize}
\item Each HTML page is a node, and each link occurrence is a directed edge.
Duplicates affect link counts and PageRank contributions, but not shortest-path
distances.
\item PageRank distributes a source page's old rank among its outgoing links
and adds the teleportation term. A page with no outgoing links distributes its
rank uniformly. All updates use the previous iteration's vector.
\item Conserving the sum of ranks does not mean that individual ranks have
stopped changing. The program exposes the literal sum-only stopping rule and
the documented additional L1 check used in the measured runs.
\item BFS finds shortest paths in a graph with unit-length edges. Outward
closeness combines the number of reachable pages with their distance sum.
The bitmask version performs the same BFS level expansion using integer sets.
\item Gzip reduces transfer size without changing the decoded HTML or graph.
Parallelism is confined to library-managed downloads; graph processing follows
after the downloads and uses one thread.
\end{itemize}
The validation uses manually specified small graphs, exact PageRank solutions,
and an independent Floyd--Warshall distance calculation. These checks help
detect errors in suggested code rather than relying only on plausible-looking
results from the generated dataset.

\clearpage
\subsection{Cloud Billing evidence and reporting delay}
The screenshot below was supplied on September 25, 2026. It shows the Billing
Reports view grouped by project, with the row for
\texttt{thermal-circle-508221-m7}. The selected setting is the current month;
the actual-cost summary covers September 1--25, 2026. The forecast for
September 1--30 is not used as actual spending.

\begin{tabularx}{\linewidth}{@{}lY@{}}
\toprule Item & Displayed value \\ \midrule
Project & \texttt{thermal-circle-508221-m7} \\
Actual-cost summary period & September 1--25, 2026 \\
Usage cost & \$0.01 \\
Other savings & \(-\$0.01\) \\
Net project subtotal & \$0.00 \\
\bottomrule
\end{tabularx}

\textbf{Scope and limitation.} These values are the displayed month-to-date
project totals, not a separately attributed final HW2 bill. This project was
also used for Homework 1, so the monthly view may contain unrelated earlier
usage. Billing had not fully updated at report preparation. The instructor
confirmed on Piazza that charges can be delayed by about 24 hours and allowed
students to demonstrate how to obtain the information when the reported values
are incomplete. The displayed \$0.00 therefore does not establish that HW2
incurred no charges; a final HW2-specific total was not yet available.

To obtain the assignment-period cost when the records update:
\begin{enumerate}
\item Select project \texttt{thermal-circle-508221-m7} in Google Cloud Console
and open \textbf{Billing $\rightarrow$ Reports}.
\item Set the custom usage-date range to September 24--25, 2026, covering the
HW2 cloud setup and execution, and restrict the project filter to this project.
\item Inspect actual costs grouped by service, including Cloud Storage and
Compute Engine. Record usage cost, savings, and net subtotal, together with
any other usage in the selected period. Do not substitute forecasted cost.
\end{enumerate}

\begin{center}
\includegraphics[width=\linewidth,height=3.95in,keepaspectratio]{\BillingImage}
\\[4pt]{\small Figure 1. Actual billing-console screenshot: current-month project
usage cost, savings, and subtotal. Charges were subject to reporting delay.}
\end{center}

\clearpage
\appendix
\section{Parameters, source files, and references}
\subsection{Command-line parameters}
\begin{tabularx}{\linewidth}{@{}p{0.32\linewidth}Y@{}}
\toprule Parameter & Meaning \\ \midrule
\texttt{--bucket NAME} & Read numbered pages from this bucket. \\
\texttt{--prefix PATH} & Object prefix; \texttt{hw2/} or \texttt{hw2-gzip/}. \\
\texttt{--local DIRECTORY} & Read local files instead of a bucket; used for local checks. \\
\texttt{--expected-nodes N} & Reject a loaded graph with a different number of nodes; use 12000. \\
\texttt{--download-workers N} & Default 1; recorded final runs use 32 library-managed download workers. Does not parallelize graph processing. \\
\texttt{--damping D} & PageRank damping factor; default 0.85. \\
\texttt{--threshold PCT} & Stopping threshold in percent; default 0.5. \\
\texttt{--convergence MODE} & \texttt{sum}: total-rank change; \texttt{l1}: vector movement; \texttt{both}: require both, the default and recorded mode. \\
\texttt{--verbose-pr} & Display each iteration's rank sum and change measurements. \\
\texttt{--authenticated} & Use application default credentials instead of anonymous access. Not required for this public bucket. \\
\texttt{--skip-closeness} & Omit closeness for troubleshooting; not used in completed measurements. \\
\bottomrule
\end{tabularx}
Exactly one of \texttt{--bucket} and \texttt{--local} is required. Parallel
download workers apply only to bucket mode. Complete runs include closeness.

\subsection{Source files and loading behavior}
\begin{itemize}
\item \texttt{generate.py}: supplied generator logic, with comments and formatting.
\item \texttt{analyze.py}: command-line interface, storage/local loading, and timings.
\item \texttt{graphlib\_hw.py}: graph representation, parsing, statistics, PageRank, and BFS.
\item \texttt{test\_pagerank.py}, \texttt{test\_closeness.py},
\texttt{test\_regressions.py}, and \texttt{test\_gcs\_io.py}: independent correctness tests and mocked storage tests.
\item \texttt{upload\_to\_gcs.py}: an alternative sequential Python uploader.
\item \texttt{results/}: measured output, environment records, package versions, tests, bucket IAM, and cleanup evidence.
\end{itemize}
The loader lists objects in pages of 100, requesting names, generations, and the
next-page token. Parallel downloads are submitted in batches of 100 into a fresh
temporary directory. Parsing starts only after all downloads succeed. The
requests use a 10-second connection timeout, a 60-second read-inactivity timeout,
and a 300-second retry budget; an in-progress request can extend beyond that
budget. Completed runs do not reuse cached downloads. Temporary files are
cleaned up on normal completion or handled exceptions.

\subsection{References}
\begin{enumerate}[label={[\arabic*]}]
\item Google Cloud, \href{https://docs.cloud.google.com/storage/docs/access-control/making-data-public}{Making data public}: bucket IAM and \texttt{allUsers}.
\item Google Cloud, \href{https://docs.cloud.google.com/python/docs/reference/storage/latest/google.cloud.storage.transfer_manager}{Python storage transfer manager}: library-managed downloads.
\item Google Cloud, \href{https://docs.cloud.google.com/sdk/gcloud/reference/storage/cp}{\texttt{gcloud storage cp}} and \href{https://docs.cloud.google.com/storage/docs/transcoding}{gzip transcoding}: content encoding and compressed objects.
\item Google Cloud, \href{https://docs.cloud.google.com/billing/docs/how-to/reports}{Cloud Billing Reports}: project/date filters and actual costs.
\item Course assignment and instructor Piazza clarifications: nine statistics per direction, one course repository, and library-managed parallel downloads with single-threaded graph processing.
\end{enumerate}
\end{document}
